\Task Дан массив $a_{1}, \dots , a_{n}, a_{i} \in \R$. Необходимо извлечь наибольшую по
длине возрастающую подпоследовательность. То есть найти такие
индексы $1 \Le i_{1} < i_{2} < . . . < i_{k} \Le n$, что $a_{i_{j}} < a_{i_{j+1}}$ и $k$ максимально за $\mathcal{O}(n\,log(n))$.

\textbf{Решение 1:}

1) Сформируем пары $\begin{pmatrix}
a_{i}\\
i
\end{pmatrix}$

2) Отсортируем их по возрастанию $a_{i}$, при равенстве - по убыванию $i$. Тогда, ограничение $a_{j} < a_{i}$ естественно выполняется.

3) \textit{Смысл dp.} dp[i] - длина НВП, заканчивающейся на $a_{i}$.

4) \textit{База dp.} dp[:] = 0.

5) \textit{Пересчёт, изменение по индексу} dp[i] = $1 + \max \limits_{j < i}(dp[j])$ - Можно посчитать с помощью \hyperlink{segtree}{дерева отрезков} за $\mathcal{O}(log(n))$.

6) \textit{Порядок вычислений. for ($a_i, i$) in sorted(pairs)} 

7) \textit{Ответ} лежит в dp[:].

\drawsomesmall{pix/picsol_nvp.jpg}

\textbf{Решение 2:}

1) \textit{Смысл dp.} dp[i] - минимальный элемент, на котором заканчивается последовательность длины $i$.

2) \textit{База dp.} dp[0] = $-\inf$, dp[1:] = $+\inf$.

3) \textit{Пересчёт} Место для $a_i$ можно найти бинпоиском. И правда, возрастающая подпоследовательность с ним может заканчиваться на это число только если оно больше последнего элемента, на которое заканчивается подпоследовательность без него, то есть оно может встать максимум на dp[k+1], где k - наибольший индекс элемента в dp, которое меньше чем $a_i$.

\begin{lstlisting}
int dp[n+1];
dp[0] = -inf;
for (int i = 1; i < n + 1; i++) {
  dp[i] = +inf;
}

for (int i; i < n; i++) {
  int left = 0;
  int right = n+1;
  while (left + 1 < right) {
    int mid = (left + right) / 2;
    if (dp[mid] < a[i]) {
      left = mid;
    } else {
      right = mid;
    }
  }

  dp[right] = a[i]
}
\end{lstlisting}
4) \textit{Порядок вычислений. for i in 0 ... n} 

5) \textit{Ответ} индекс последнего элемента, который лежит в dp, меньший $+inf$.
\pagebreak